\section{Same-task error composition and paid implementation}\label{sec:errors}
First-discrepancy coupling, state-product simulation, and attained hidden-offset and erased-bit constructions are inherited mechanisms~\cite{GTFStopped,GTFBayes}. We state their precise form for the present mean-budget architecture, and combine them with its correlated acquired-geometry law.
\subsection{Stopped-law perturbations}
A numerical or calibration comparison must include the stopping decision and the protected audit. Approximation of the report marginal alone does not compare the scored task.

\begin{proposition}[Acquired-occupation error composition]\label{prop:stopped-error}
Consider two causal implementations on the same visible interface. Suppose their next conditional interface kernels, on a common pre-acquisition history $h$, admit a coupling with mismatch probability at most $\kappa(h)$; the protected final audit coupling has mismatch at most $\kappa_f(h)$. Use the same observer and fresh randomization in both. If the ideal implementation has finite mean stopping time, then the total variation distance of their complete stopped, scored laws is at most
\begin{equation}\label{eq:stopped-tv}
 \min\left\{1,\ \E\sum_{t=1}^{\tau}\kappa(H_{t-1})
                    +\E\kappa_f(H_\tau)\right\}.
\end{equation}
A possible nonterminating approximate run is recorded at an extra failure outcome. If $0\le\ell\le L$, the corresponding risk difference is at most $L$ times \eqref{eq:stopped-tv}. In particular, uniform bounds $\kappa\le\theta$, $\kappa_f\le\theta_f$ and $\E\tau\le N$ give $L(N\theta+\theta_f)$, truncated at $L$.
\end{proposition}
\begin{proof}
Couple the conditional kernels sequentially while the histories agree, and use identical routing coins. The event that the first disagreement occurs at acquisition $t$ has probability at most the ideal expectation of $\one_{\{\tau\ge t\}}\kappa(H_{t-1})$. After the ideal stop, couple the protected audit in the same way. The union bound and monotone convergence give \eqref{eq:stopped-tv}; on its complement the complete finite stopped records and audit outcomes agree. Coupling bounds total variation with convention $\norm{P-Q}_\TV=\sup_A|P(A)-Q(A)|$. Integrating a loss in $[0,L]$ gives the risk bound. This argument concerns the conditional kernels of the full compared interface, including the hidden-state effect needed for later reports or its equivalent visible conditional law. A bound on a single current-report marginal is not substituted for that hypothesis.
\end{proof}

\begin{remark}[Mean cost is a separate obligation]\label{rem:mean-not-tv}
Total variation does not control expected acquisition. Replacing a one-call stop by a branch of probability $\eps$ lasting $\eps^{-2}$ calls changes the complete-law probability of that branch by at most $\eps$, while adding order $\eps^{-1}$ to the mean. Thus \eqref{eq:stopped-tv} is not a budget certificate for the approximate machine. For the rejection construction one may plan an ideal conditional stop floor $(1+\zeta)/N$ and require the implemented acceptance probability to differ by at most $\zeta/N$ at every surviving history. The implemented floor is then $1/N$, which proves its own mean bound and geometric tail. Appendix~\ref{app:computation} gives a geometric alternative with a deterministic acceptance margin.
\end{remark}

If a fixed acceptance rule is retained but each final code vector is perturbed by at most $\delta_{\rm out}$, Minkowski's inequality gives
\begin{equation}\label{eq:root-error}
 \E\norm{X_\tau-\widehat r_J}^2
   \le\left(\sqrt{\E\norm{X_\tau-r_J}^2}+\delta_{\rm out}\right)^2.
\end{equation}
Only when $r_J=\E[X_\tau\mid J]$ under the same actual stopped law does orthogonality improve this to addition of $\delta_{\rm out}^2$. Centers optimized for a reference law need not be calibrated centroids for every unknown correlated model. The stronger additive formula is therefore not asserted for them.

\subsection{Resource-certified causal simulation}
\begin{proposition}[Complete-task causal transfer]\label{prop:morphism}
A parameter-independent causal simulator from a source experiment to a target experiment has at most $R$ internal states and a causal action lift. It transforms reports, stopping, and the protected task readout compatibly. Suppose that, uniformly over models and all target policies of interest, its complete stopped scored law has total variation error at most $\eps$, and its source acquisition certificate is
\[
 \E[\text{source calls}]\le\kappa\E[\text{target calls}]+c.
\]
Then a target observer with $M$ states and mean acquisition at most $N$ has a source implementation with at most $RM$ states and mean acquisition at most $\kappa N+c$, satisfying
\begin{equation}\label{eq:morphism-risk}
 \calR_{\rm source}(\kappa N+c,RM)
       \le\calR_{\rm target}(N,M)+L\eps.
\end{equation}
Two such simulators compose with error at most $\eps_1+\eps_2$, state multiplier $R_1R_2$, and acquisition bound
\begin{equation}\label{eq:morphism-cost}
 N\longmapsto\kappa_1\kappa_2N+\kappa_1c_2+c_1.
\end{equation}
Program descriptions, temporary workspaces and physical execution times remain separately charged.
\end{proposition}
\begin{proof}
Run the simulator and target observer in their product state. Only past reports and fresh coins enter the product transition. The action lift ensures that every requested target action has a legal source realization; compatibility of the stop and audit prevents information from arriving after the target decision but before the source decision. Apply the declared mean-cost certificate to the executed target policy and the total variation loss bound to the complete task. Infimizing over target observers proves \eqref{eq:morphism-risk}. Sequential composition gives \eqref{eq:morphism-cost} by substitution. Kernel contraction of total variation and the triangle inequality give the additive error. The product of internal state sets gives the state count. Reverse lower bounds require a reverse simulator or an independent lower argument; they do not follow from this one-sided construction.
\end{proof}

\subsection{An attained joint resource--risk law}
It is useful to distinguish an allowable error upper bound from an experiment that actually has that error. We give the latter in the same task. Start with a regular-support report-recoverable experiment from Theorem~\ref{thm:main}. A stopped core state is frozen. The protocol then requires exactly one bit acquisition: for an unknown model bit $b\in\{-1,1\}$ it reports $b$ with probability $1-e$ and an erasure symbol with probability $e$. The bit affects no preceding core report. An unknown calibration parameter $c\in\{-\delta,\delta\}$ affects no observer report at all. The final protected audit returns the core audit, $b$, and $c$, and the loss is the sum of their squared prediction losses. The observer must commit all three predictions before this audit. It is forbidden to restart the core or repeat the bit acquisition.

\begin{lemma}[Attained erasure deficiency]\label{lem:erasure}
For this bit interface, relative to the perfect bit experiment, the minimum parameter-independent causal simulation error is exactly $\eps=e/2$. The protected truth readout is fixed and cannot be altered by the simulator. The minimax bit prediction risk is $e$, and the independent hidden-calibration prediction risk is $\delta^2$.
\end{lemma}
\begin{proof}
On a nonerased report, output the reported bit. On erasure, a fair guess gives error $e/2$ for each model bit, so it achieves the stated total variation defect against the perfect report and unchanged protected truth. Conversely, under the uniform prior on $b=\pm1$, no pre-audit information on an erased branch depends on $b$. Every simulator then guesses wrongly with probability at least $1/2$ on that branch. Its largest model error is at least its prior average, $e/2$. This proves exact deficiency, without access to the final audit. For squared bit prediction, the conditional Bayes estimator on erasure is zero and its loss is one; the minimax risk is $e$, attained by zero on erasure and the reported sign otherwise. Since the calibration sign is absent from every report, averaging its two signs gives $\E(r-c)^2=\E r^2+\delta^2$, minimized at zero.
\end{proof}

\begin{theorem}[Attained acquisition--memory--error composition]\label{thm:joint-law}
Let $0\le\delta\le1$, $0\le\eps\le1/2$, $n\ge1$ and $M\ge5$. Give the core a mean acquisition budget $n$ and use the extended protocol with erasure probability $2\eps$. Its model-uniform optimum satisfies
\begin{equation}\label{eq:joint-law}
 \calR_{\rm ext}(n,M,\delta,\eps)-B
        \asymp (nM)^{-2/s}+\delta^2+\eps.
\end{equation}
The upper is implemented by at most $4m+1\le M$ states, with $m=\lfloor(M-1)/4\rfloor$. Under the convention charging one preparation, each core call, one bit call and one protected audit call, the mean total physical-call budget is $\mathsf N=n+3$. Thus the geometric term in that convention is $((\mathsf N-3)M)^{-2/s}$, not $(\mathsf NM)^{-2/s}$ with the extra calls silently omitted.
\end{theorem}
\begin{proof}
For the upper, use the one-live-state core construction with $m$ centers. Upon acceptance, retain its center index in one of $m$ bit-waiting states. These states perform the single bit call. Its three possible reports lead to $3m$ terminal states, carrying the fixed core center and the bit prediction in $\{-1,0,1\}$. The calibration prediction is zero. The total is $1+m+3m=4m+1$. No clock, free retained bit or exact posterior is present. The core risk bound and Lemma~\ref{lem:erasure} give
\[
 \calR_{\rm ext}-B\le C(nm)^{-2/s}+\delta^2+2\eps.
\]
Since $m\ge M/8$ for $M\ge5$, this is the upper in \eqref{eq:joint-law}.

For the converse, any legal extended observer commits to a core stop before receiving the bit. Its actual stopped core law is bounded by $b_*n\nu$. The first component of its final decoder has at most $M$ possible values, so its core squared distortion is at least $\calQ_M^\nu(b_*n)\ge c(nM)^{-2/s}$. This remains true even though the later bit report can select among those values. Average over the two bit signs and two calibration signs, keeping any one core model fixed. These signs do not change the preceding core law. Conditional on an erased bit there is still no information about its sign; there is never pre-audit information about the calibration sign. The three squared losses add, so their averaged lower bound is $c(nM)^{-2/s}+2\eps+\delta^2$. The worst model risk is at least this average. The physical-call count is a direct count of the specified protocol. There is no conversion from total variation to an expected-cost bound in this proof.
\end{proof}

Theorem~\ref{thm:joint-law} supplies attained lower and upper error terms for a fixed enlarged task. Proposition~\ref{prop:morphism} is the general transfer statement. They are logically different results: an arbitrary permitted calibration tolerance or simulator error need not be attained by a given experiment and hence need not cause that much risk.

\subsection*{Same-budget primitive perturbations}
For common ordered carriers, Lemma~\ref{lem:resolvent-error} strengthens the bounded-task comparison by a separate occupation-weighted mean-cost estimate after program-side stopping regularization. Proposition~\ref{prop:same-budget} then transfers the optimum at the same $N,M$ under uniform raw instrument and preparation defects. It does not assert that total variation of an arbitrary stopped transcript controls its expected length. A simulator that changes the carrier, consumes extra native calls, or retains additional state must still satisfy Proposition~\ref{prop:morphism}; the stopping argument does not erase its cost.
